\section{Conclusion}
\label{s:conclusion}

This paper investigates whether the Brasília subway raised income in the areas it serves and, in particular, whether that growth reflects real income gains among the residents who already lived there rather than a shift in who lives there. Using census-tract data for 2000 and 2010 and an instrumental-variable strategy based on a planned and discarded subway alignment, we find that tracts exposed to the subway experienced income growth approximately 10 to 11 percent higher than comparable unexposed tracts over the decade, an estimate that remains positive and of similar magnitude across sample radii, distance thresholds, continuous measures of exposure, and alternative standard errors. We then test whether this gain could come from something other than higher incomes among incumbent residents. Subway exposure does not raise population, households, apartment share, or urban land cover (population and household growth are, if anything, lower in exposed tracts), and it produces no coherent shift in the demographic composition of exposed tracts, so neither physical restructuring nor the replacement of lower-income residents by wealthier ones explains the result. In addition, relative to comparable unexposed tracts, the share of household heads with positive income rises by about 1.7 percentage points and the average income of those who already had income rises by about 9 percent, consistent with the labor-market channel. To summarize, the evidence indicates that the income gains from the subway reflect real income growth among the residents who already lived in the areas it serves, and not a change in the physical structure of these areas or in who lives there.

Brasília's urban configuration makes it hard for the residents of its satellite cities to share in the income gains of the city's productive core. The \textit{Plano Piloto} region alone holds nearly half of the jobs in the Federal District (47.72 percent in 2011; \citealp{miragaya2013perfil}) and is legally protected from densification, so these residents cannot move closer to jobs, and the gains are more likely to be capitalized into housing prices than to reach them. In the areas served by the subway, however, the results point to higher incomes among the residents who already lived there, with no sign of physical restructuring or of newcomers replacing incumbents; the gains thus reached people who could not move closer to jobs. The network today has 27 stations along 42 km of track and serves six administrative regions \citep{institutomdt2019evolucao}. Thus, we understand that extending mass transit to peripheral regions the network does not yet reach could help raise the incomes of the residents who already live there, without depending on changes to the land-use rules of the \textit{Plano Piloto}.

Finally, Brasília's experience offers lessons beyond its local context. In many cities, constraints on housing supply keep workers from moving closer to jobs \citep{glaeser2018economic}, and Brasília is an extreme version of that situation. The introduction asked whether, when moving closer to jobs is not an option, transit can still raise the incomes of the workers who stay. Our estimates for Brasília, a developing-country city where evidence on this question is thinner, indicate that it can. Unlike earlier studies of the United States and France, which attribute part or all of the local gains near new stations to changes in who lives or works in the area \citep{heilmann2018transit, mayertrevien2017largescale, tyndall2021local}, the gains documented here appear to reflect income growth among residents who were already there. Because the estimates come from one city and one decade, they show that mass transit can raise the incomes of incumbent residents in a setting like Brasília's, not that it will do so everywhere. Even so, for cities that cannot densify their centers, Brasília's experience suggests that new transit lines can benefit the people who stay in place, and not only those who can afford to move in.

% \section*{Appendix}
% \addcontentsline{toc}{section}{Appendix}
